\documentclass[11pt]{article}
\usepackage[margin=1in]{geometry}
\usepackage{amsmath,amssymb}
\usepackage[hidelinks]{hyperref}
\setlength{\parindent}{0pt}
\setlength{\parskip}{0.7em}

\title{A literature answer to the general surface-area-measure question\\
for log-concave functions}
\author{Status note for arXiv:2006.16933}
\date{19 July 2026}

\begin{document}
\maketitle

\section*{Status}

The open problem in L.~Rotem, \emph{Surface area measures of log-concave
functions}, arXiv:2006.16933, is answered by the same author's later paper
\emph{The anisotropic total variation and surface area measures},
arXiv:2206.13146.  The answer is Definition~1.2, Definition~1.4, and
Theorem~1.5 of the later paper.

\section*{The source question}

For an integrable upper-semicontinuous log-concave function
$f=e^{-\varphi}$ on $\mathbb R^n$, arXiv:2006.16933 uses the push-forward
\[
  S_f=(\nabla\varphi)_{\#}(f\,dx).
\]
This correctly gives the interior contribution for sufficiently regular
functions, but if $f=\mathbf 1_K$ is the indicator of a convex body, it reduces
to $|K|\delta_0$ and loses the classical surface-area measure $S_K$.
Accordingly, PDF page~5 asks for an extension of its Definition~1.4 and
first-variation Theorem~1.5 that works for every log-concave $f$, including
both smooth functions and indicators of convex bodies.

\section*{The later answer}

Let $K_f=\overline{\{x:f(x)>0\}}$.  ArXiv:2206.13146, Definition~1.2,
associates to $f=e^{-\varphi}$ two measures:
\[
  \mu_f=(\nabla\varphi)_{\#}(f\,dx)
  \quad\text{on }\mathbb R^n,
  \qquad
  \nu_f=(n_{K_f})_{\#}
  \left(f\,d\mathcal H^{n-1}\big|_{\partial K_f}\right)
  \quad\text{on }S^{n-1}.
\]
Here $n_{K_f}$ is the almost-everywhere defined outer Gauss map.  Definition
1.4 calls the pair $(\mu_f,\nu_f)$ the surface-area measures of $f$.

Theorem~1.5 on supporting PDF page~3 proves, for arbitrary
$f,g\in\mathrm{LC}_n$ with $0<\int f<\infty$, the unrestricted formula
\[
  \delta(f,g)
  =\int_{\mathbb R^n}h_g\,d\mu_f
   +\int_{S^{n-1}}h_{K_g}\,d\nu_f,
\]
with no smoothness or strict-convexity assumptions.  The paper further
explains on PDF pages~4--5 that the indicator perturbation $g=\mathbf 1_L$
is an anisotropic coarea formula.

For $f=\mathbf 1_K$, the definitions give
\[
  \mu_f=|K|\delta_0,
  \qquad
  \nu_f=S_K.
\]
Thus the second term restores exactly the boundary contribution missing from
the one-measure definition, while for full-support smooth functions the
boundary is empty and $\nu_f=0$.  This is precisely the requested common
extension.

\section*{Conclusion}

The question is not a viable new-proof target: arXiv:2206.13146 supplies the
general definition and the complete first-variation theorem, with the
indicator-of-a-convex-body case built into the same framework.

\section*{References}

L.~Rotem, \emph{Surface area measures of log-concave functions},
arXiv:2006.16933 (2020), especially PDF page~5.

L.~Rotem, \emph{The anisotropic total variation and surface area measures},
arXiv:2206.13146 (2022), Definitions~1.2 and~1.4 and Theorem~1.5,
PDF pages~2--5.

\end{document}
